% Research-agenda addition (attribution strand): which part of a skill helped.
% Inserted after the closing pointer of the skill-helpfulness subsection.
\subsection{Which part of a skill helped}\label{sec:stp-skill-attribution}

Protocol~\ref{proto:skill} records that a skill entered the context. A skill
is a bundle: instructions in a file called \texttt{SKILL.md}, reference
files, scripts and examples. The agent reads some of them and not others. In
one documented loader, a skill's name and description sit in the system
prompt from the start; the instructions enter the context when the skill
triggers; other files enter only when read; and a script's code never enters
the context, only its output~\pdcite{anthropic2026skills}. A run therefore
shows an \pdterm{exposure ladder}, the ordered levels of contact an agent can
have with a skill: listed, loaded, file read, script run.

Never use ``opened'' as the treatment. The agent decides to open a file
after seeing the task and its own progress, so a comparison of opened with
unopened runs mixes the file's effect with whatever made the agent open it.
Not opening is not the absence of an effect either: a skill's listed name
and description can change behaviour when the skill is never
invoked~\pdcite{zhang2026estimands}.

A trace cannot settle the question. A study that paired each skill run with
a no-skill or matched-skill reference run on the same task confirmed $307$
skill-induced failures. Of these, $182$ were efficiency regressions: both
runs passed, but the skill run used more tokens and more time, at least
twice as much on one of the two. The largest sources of procedure the skill
imposed were excessive verification ($67$ cases) and heavy implementation
pipelines ($30$)~\pdcite{dong2026harmful}.
In such a run the skill's files were opened and its procedure followed; the
trace looks like diligent use. Only the reference run shows that the skill
induced the cost.

\begin{protocol}[Masked-asset extension of the skill event]\label{proto:stp-masked-asset}
\pdstatus{\Designed}{not implemented in the reference system; no result from it is reported here}
Extend each event of Protocol~\ref{proto:skill} as follows.
\begin{enumerate}[leftmargin=1.5em,itemsep=1pt]
  \item \textbf{Exposure ladder.} Record listed, loaded, each file read
  (path and byte range), each script run (exit status and an output hash),
  and the tokens of skill content in the context.
  \item \textbf{Harbor-drawn assignment.} The harbor, not the skill, draws
  whether the skill is available, unavailable, or listed only. It logs the
  probability of each arm, and every eligible arm keeps a probability of at
  least $\varepsilon>0$, so that logged runs support doubly robust off-policy
  estimates~\pdcite{dudik2011dr}.
  \item \textbf{Asset mask.} On an audit share of runs, each asset is hidden
  independently with probability $\tfrac12$. Hiding removes the file
  \emph{and} its mention in \texttt{SKILL.md}, so a dangling reference does
  not itself change behaviour. What the mask estimates is therefore the effect
  of the file and its mention together, not of reading the file.
  \item \textbf{Commitment.} Before the run starts, the draws are committed
  to the session's \pdterm{evidence trail}, its sequence of timestamped notes,
  to which notes can be added but never changed.
\end{enumerate}
\end{protocol}

\begin{table}[H]
\centering\small
\caption{What the masked-asset protocol can estimate, and what each estimate supports.}
\label{tab:stp-skill-estimands}
\begin{tabularx}{\textwidth}{@{}l >{\raggedright\arraybackslash}X >{\raggedright\arraybackslash}X@{}}
\toprule
\textbf{Estimand} & \textbf{Contrast} & \textbf{Supports} \\
\midrule
Availability effect & skill available vs.\ unavailable, all tasks & deranking and settlement \\
Asset main effects & each asset shown vs.\ masked, audit share & naming a helpful asset, if interactions are small \\
Effect of reading & the mask as an instrument for reading: its effect on outcome divided by its effect on the read rate & not supported here: the mask also removes the mention, which can act without any read, so report the joint file-and-mention effect \\
Edit effect & two content hashes, randomized & ``this edit improved the skill'' \\
\bottomrule
\end{tabularx}
\end{table}

\keyidea[What the system prompt is responsible for]{The same question arises
one level up: how much of an outcome is due to the system prompt, rather
than the model, the task or the skills? A crossed design with variance
components answers a narrower question, namely how much variance each factor
carries in the levels sampled. Lewontin showed that such an analysis is ``a
unique function of the present distribution'' of causes, and that when
causes interact ``it becomes conceptually impossible to assign quantitative
values to the causes of that individual event''~\pdcite{lewontin1974}.

Interaction is the rule here. Meaning-preserving format changes moved
few-shot accuracy by up to $76$ points for LLaMA-2-13B, and format
performance correlated only weakly between models~\pdcite{sclar2024}. In a variance decomposition
of open agent benchmarks, the agent main effect was below $3\%$ of total variance in every
cell, while the agent-by-task interaction was $7.5$ to
$12.5\%$~\pdcite{srinivasan2026ddr}. A Shapley split over factors does not
resolve this: it divides each interaction effect evenly among the factors
involved. That is a convention, not a measurement.

Report the prompt's effect as a contrast: under a stated model, harness and
task distribution, changing $P$ to $P'$ moves the outcome by $\Delta$, with
its interval. Add the spread across meaning-preserving variants of $P$ and
the share of variance in prompt-by-model and prompt-by-task interaction.
Hash the system prompt with every run, beside the model and skill versions
of Table~\ref{tab:stp-not-bandit}.}

\noindent Settlement on these estimates is the subject of
\pdchapref{he}{The Harbor Economy}%
\ifpdbook
\ (\S\ref{he:sec:he-licence-split})%
\fi
.

\pdexercisepointer{\ref{ex:stp-opened-not-treatment}}{\ref{ex:stp-prompt-contrast}}{\pageref{ex:stp-opened-not-treatment}}
